\section{实验评估}
\label{sw-eval-section}

\paragraph{材料与执行。}
实验围绕技能使用质量、证据治理和来源变化后的演化展开，采用Petstore/OpenAPI的创建、查询和删除任务。
材料包含60个任务输入、30对共60个治理状态，以及12项路径、请求必填项和允许值变化；每项变化包含8个目标题和8个保持题。
模型为DeepSeek-V4.1-Flash-1017（API别名deepseek-flash），temperature为0，自动重试为0，任务与提案输出上限分别为2048和4096 tokens。
主路线使用SkillWiki本地Harness，评价请求构造和已给定响应的解释，不联网调用业务服务。
归档记录包含4590次唯一模型调用，其中18次为更新提案；技术重复与独立材料数量分别统计。

\paragraph{技能使用质量。}
同一输入在无领域材料、直接提供来源、提供固定技能三种条件下各重复两次。
三种条件分别答对84/120、112/120、120/120，全部执行完成且格式符合要求。
技能条件比直接使用来源提高6.67个百分点，创建、查询和删除操作各为40/40。
此处模型原答与最终输出相同，说明所选技能表示能够帮助模型更准确地使用来源知识。

\paragraph{证据治理。}
治理状态覆盖正常发布、身份与内容、适用性、完整性、暂缓解除、恢复与状态一致性。
独立Git/CI+对照实现相同证据合同并独立执行仓储动作。
两路线在主实验各50/50、补充修订各10/10个状态上决定正确，误放行及状态不一致均为0。
主实验中SkillWiki有1项入口验证错误符合预登记允许决定集合，数据保留原分类。
在所测证据要求上，SkillWiki原生治理与认真编写的Git/CI+实现一致，体现了这些Git/CI治理能力在技能生命周期中的整合。

\paragraph{来源更新与生命周期。}
12项变化比较原规则参照、未更新技能和更新技能，每输入重复两次，共1152个核心回答。
按新规则评价，目标回答由未更新条件的0/192提高到更新后的原答176/192、最终192/192，其中16次由显式来源校验纠正。
原有能力保持为173/192，未更新条件为192/192；7项事件满足改善目标且无保持退步的联合判据。
每身份另需完整36题资格验证，主实验两路线各有5项更新合格并完成发布、恢复后重验及再发布。

在请求与响应约束分离的后续修订中，五项修订候选最终资格为180/180、原答171/180；
五个恢复新身份最终180/180、原答167/180，均拒绝旧证明，并在重验后发布和回读。
该候选集合包含四个R1身份和一个R2身份，U06的R1失败35/36仍予保留。
U08补充验证的更新与恢复身份原答及最终均为36/36，完成发布与回读。
跨记录修订，11项事件各有适用版本的成功链；U03仍为35/36，越界数值虽被拒绝，但错误分类不符合合同，因此未发布。

\paragraph{结果与材料范围。}
实验展示了所选技能的使用价值、与独立Git/CI+一致的证据决策，以及验证约束下的更新和恢复闭环。
材料属于同一API家族的开发、回归与相关扩展；各结果保留运行与候选版本，不合为同一版本的整体通过率。
公开配套材料提供合成输入、预期答案、来源快照、治理状态、结果表和脱敏调用索引。
